%================================================================================
%       Safety Critical Systems Club - Data Safety Initiative Working Group
%================================================================================
%                       DDDD    SSSS  IIIII  W   W   GGGG
%                       D   D  S        I    W   W  G   
%                       D   D   SSS     I    W W W  G  GG
%                       D   D      S    I    WW WW  G   G
%                       DDDD   SSSS   IIIII  W   W   GGG
%================================================================================
%               Data Safety Guidance Document - LaTeX Source File
%================================================================================
%
% Description:
%   Miscellanea section.
%
%================================================================================
\chapter{Data Properties\index{Data Property|(}}\label{bkm:dataProperties}

\dsiwgSectionQuote{Failure is an amazing data point that tells you which direction not to go.}{Payal Kadakia}

\Glspl{data property} are used to establish what aspects of the data (e.g., \gls{timeliness}\index{Timeliness Property}, \gls{accuracy}\index{Data Property!Accuracy}\index{Accuracy Property}) need to be guaranteed in order that the system operates in a safe manner.

\autoref{tab:PropertiesOfData} documents a non-exhaustive collection of \glspl{data property}. Typically, it is the loss of one of these \glspl{data property} that presents a \gls{hazard}. This notion of ``loss'' is dependent on the intended use; for example, what is ``timely'' for one use may not be for another.

\begin{longtable}{|L{\dsiwgColumnWidth{0.24}}|C{\dsiwgColumnWidth{0.15}}|L{\dsiwgColumnWidth{0.61}}|}
	\caption{Properties of data}
	\label{tab:PropertiesOfData}
	\\\hline\TableHeadColour{Property} & \TableHeadColour{Abbreviation} & \TableHeadColour{Description}\\\hline
	\endfirsthead
	\caption[]{Properties of data (continued)}
	\\\hline\TableHeadColour{Property} & \TableHeadColour{Abbreviation} & \TableHeadColour{Description}\\\hline
	\endhead
	\multicolumn{3}{r}{\sl Continued on next page}
	\endfoot
	\endlastfoot
	\Gls{integrity}\index{Data Property!Integrity}\index{Integrity!Property} & I & The data is correct, true and unaltered\\\hline
	\Gls{completeness}\index{Data Property!Completeness}\index{Completeness!Property} & C & The data has nothing missing or lost\\\hline
	\Gls{consistency}\index{Data Property!Consistency}\index{Consistency!Property} & N & The data adheres to a common world view (e.g., units)\\\hline
	{\Gls{continuity}}\index{Data Property!Continuity}\index{Continuity Property} & Y & {The data is continuous and regular without gaps or breaks}\\\hline
	{\Gls{format}}\index{Data Property!Format}\index{Format Property} & O & {The data is represented in a way which is readable by those that need to use it}\\\hline
	{\Gls{accuracy}}\index{Data Property!Accuracy}\index{Accuracy Property} & A & {The data has sufficient detail for its intended use}\\\hline
	{\Gls{resolution}}\index{Data Property!Resolution}\index{Resolution Property} & R & {The smallest difference between two adjacent values that can be represented in a data storage, display or transfer system}\\\hline
	{\Gls{traceability}}\index{Data Property!Traceability}\index{Traceability Property} & T & {The data can be linked back to its source or derivation}\\\hline
	{\Gls{timeliness}}\index{Data Property!Timeliness}\index{Timeliness Property} & M & {The data is as up to date as required}\\\hline
	{\Gls{verifiability}}\index{Data Property!Verifiability}\index{Verifiability Property} & V & {The data can be checked and its properties demonstrated to be correct}\\\hline
	{\Gls{availability}}\index{Data Property!Availability}\index{Availability Property} & L & {The data is accessible and usable when an authorized entity demands access}\\\hline
	{\Gls{fidelity}}\index{Data Property!Fidelity}\index{Fidelity Property} & F & {How well the data maps to the real-world entity it is trying to model}\\\hline
	{\Gls{priority}}\index{Data Property!Priority}\index{Priority Property} & P & {The data is presented / transmitted / made available in the order required}\\\hline
	{\Gls{sequencing}}\index{Data Property!Sequencing}\index{Sequencing Property} & Q & {The data is preserved in the order required}\\\hline
	{\Gls{intended_destination}}\index{Data Property!Intended Destination}\index{Intended Destination Property} & U & {The data is only sent to those that should have access to it}\\\hline
	{\Gls{accessibility}}\index{Data Property!Accessibility}\index{Accessibility Property} & B & {The data is visible only to those that should see it}\\\hline
	{\Gls{suppression}}\index{Data Property!Suppression}\index{Suppression Property} & S & {The data is intended never to be used again}\\\hline
	{\Gls{history}}\index{Data Property!History}\index{History Property} & H & {The data has an audit trail of changes}\\\hline
	{\Gls{lifetime}}\index{Data Property!Lifetime}\index{Lifetime Property} & E & {When does the safety-related data expire}\\\hline
	{\Gls{disposability}}\index{Data Property!Disposability}\index{Disposability Property} & D & {The data can be permanently removed when required}\\\hline
	{\Gls{goldilocks}}\index{Data Property!Goldilocks}\index{Goldilocks Property} & G & {The data is just the right size -- not too much and not too little}\\\hline
	{\Gls{analysability}}\index{Data Property!Analysability}\index{Analysability Property} & Z & {The data (including any \gls{metadata}) is of a suitable size, type and format to enable it be usefully analysed}\\\hline
	{\Gls{explainability}}\index{Data Property!Explainability}\index{Explainability Property} & X & {The data can be meaningfully explained,by a suitable mechanism, to those who need to understand it}\\\hline
	{\Gls{clarity}}\index{Data Property!Clarity}\index{Clarity Property} & K & {The data is visible, unmasked, not hidden or obscured and not camouflaged}\\\hline
\end{longtable}

\section{Recently Added Properties}
\subsection{The \glsfmttext{goldilocks} Property}
The ``\Gls{goldilocks}''\index{Goldilocks Property} property\footnote{The authors are aware that some of the names of categories are very culturally dependent. For example, we are informed that the folk-tale of Goldilocks is not well known in Chinese culture but the idiom gu\`oy\'oub\`uj\'\i, interpreted as ``doing too much is just as bad as not doing enough'' is well known. The intention of the names is simply to evoke the associations, so users of this guidance are encouraged to use names for the categories that are more appropriate to their cultures.} addresses appropriate sizing and quantity of data. A number of issues have been found to arise when there is too much or too little data. While it is particularly relevant to communications links, it may have relevance to other areas such as \glspl{database} and when people are involved in reviewing or checking data. The property is named ``\gls{goldilocks}'' as it refers to the need to have not too much, not too little, but just the right amount of data%
%
\footnote{A system where the property was lost involved a high-speed data bus that connected several safety-critical systems. A transceiver of that bus failed and transmitted random noise. The receivers employed parity checks and cyclic redundancy check, but the system had been designed to eliminate occasional \glspl{error}. When random noise filled the bus, several apparently valid messages were created every second, resulting in potentially lethal behaviour.}	
\footnote{In a \gls{hazop}\index{HAZOP} carried out during 2020, based upon the \gls{hazop}\index{HAZOP} guidewords in this guidance, the facilitator realised that certain failure modes had not been identified by the \gls{hazop} team. In addition to the issue of system overload already discussed, those omissions also concerned system behaviour following data rejection. In these cases, bad data was detected and rejected, but the consequences of data rejection over an extended period had not been considered.}

The \gls{goldilocks}\index{Goldilocks Property}\index{Data Property!Goldilocks} property is related to the data volume problem mentioned in \dsiwgExtRef{Chapter}{2}{bkm:additionalways}, however, given the importance of data sizing and the experience of real-world incidents this is now a separate property.
%
\subsection{The \glsfmttext{analysability}\index{Analysability Property} property}\label{bkm:guidance:analysability}
The \gls{analysability} property recognizes that data is now highly complex and extensive, often large scale, and distributed. And yet, for safety purposes we need to make sure it is of suitable quality and able to support system goals.
We therefore need to ensure that it is possible to analyse it for key characteristics and establish meaningful results using tools or other means. This property is related to explainabilty and may be performed by the same set of tools or techniques.

\subsection{The \glsfmttext{explainability}\index{Explainability Property} property}\label{bkm:guidance:explainability}
\Gls{explainability} describes the ability to establish what the purpose and effect data (and especially changes to data) has on a system
and explain this to relevant \glspl{stakeholder} in terms that they can understand. It is particularly important for learning and AI-based systems.
This could be, for example, \gls{ml} training\index{Training!AI} data or system \gls{configuration data}.
This property is related to \gls{analysability}\index{Analysability Property} and may be performed by the same set of tools or techniques.

\subsection{The \glsfmttext{clarity}\index{Clarity Property} property}
The clarity property recognises the case where although data might be available it might not be clear to those who need it that it \textit{is} available or where it can be found.
%

\index{Data Property|)}
